\begin{statement}[C.7.5; Proved computationally: the specified integer-relation tests]
At all eleven rational CM parameters, an independent theta
calculation and the global height calculation agree to more
than $245$ decimal places. At $200$ decimal digits, PSLQ
with tolerance $10^{-185}$, maximum coefficient $10^8$ and
$10000$ steps returns no relation in each prescribed list
consisting of $\log N$, logarithms of the bad primes together
with $\log2,\log3$, $\pi$, and independent real period
components. At $j=287496$ the known identity
$\omega_1=2|\Im\omega_2|$ is removed first; repeated prime
logarithms and zero period components are likewise omitted.
An additional direct algebraicity search tests both $N$ and $N^{12}$
at all eleven points, using the vectors $(1,x,\ldots,x^d)$ for
$d=4,8,12$, working precision $235$ decimal digits, tolerance
$10^{-210}$, maximum coefficient $10^{12}$ and $20000$ steps.
All $66$ calls return no relation.
\end{statement}
\begin{proof}[Computational verification]
The complete inputs, $255$-digit values, comparison errors,
PSLQ labels and returns are in
\path{code/C_cm_250.gp}, \path{code/C_cm_250.out},
\path{code/C_cm_theta_pslq.py}, and its output. The theta
calculation uses the exact CM period ratios and solves the
Fourier equation for $\wp(1/2-iy)=0$ independently of PARI's
height implementation. It verifies the singular modulus
as well. The additional polynomial search is reproduced by
\path{code/C7_norm_polynomial_search.py}; its output and JSON
record include every parameter, source hash and returned vector.
These tests prove no transcendence assertion. A null PSLQ return
is not an exhaustive certificate excluding all polynomials within
the stated coefficient and degree limits.
\end{proof}
